(a) Assume first that $\mathfrak m\ne0$. Choose a transcendence basis $t_1,\ldots,t_r$ of $L/K$. The local domain $\mathcal O[t_1,\ldots,t_r]_{\mathfrak m\mathcal O[t_1,\ldots,t_r]}$ dominates $\mathcal O$ and has fraction field $K(t_1,\ldots,t_r)$. Thus we may assume $L/K$ finite.

Choose a valuation ring dominating $\mathcal O$. Among generators $x_1,\ldots,x_n$ of $\mathfrak m$, choose $x_1$ of least value. Then $\mathcal O'=\mathcal O[x_2/x_1,\ldots,x_n/x_1]$ lies in that valuation ring, and $x_1$ remains a nonunit. For a prime $\mathfrak p$ minimal over $(x_1)$, the principal ideal theorem gives $\dim\mathcal O'_{\mathfrak p}=1$. Since $\mathfrak m\mathcal O'=(x_1)$, its contraction to $\mathcal O$ is $\mathfrak m$, so $\mathcal O'_{\mathfrak p}$ dominates $\mathcal O$.

By Krull--Akizuki, the integral closure of this one-dimensional noetherian local domain in $L$ is noetherian and one-dimensional. Localizing at a maximal ideal over its maximal ideal gives a one-dimensional noetherian integrally closed local domain, hence a DVR dominating $\mathcal O$.

If $\mathcal O=K$ and $\operatorname{trdeg}_KL>0$, start instead with $K[t_1,\ldots,t_r]_{(t_1)}$ and apply the finite-extension step. If $\mathcal O=K$ and $L/K$ is finite, the assertion has an exception: every valuation ring of $L$ containing $K$ contains all of $L$, since its elements are integral over $K$ and valuation rings are integrally closed. Thus there is no nontrivial DVR in this case.

(b) Necessity follows from the valuative criteria. For separatedness, let $z_1\in\Delta(X)$ specialize to $z_0$ in $X\times_YX$. If the points differ, apply (a) to the local ring of $z_0$ in the reduced closure of $z_1$, with fraction field $k(z_1)$. This gives a DVR and a map of its spectrum to $X\times_YX$ taking the generic and closed points to $z_1,z_0$. The two projections agree at the generic point, so uniqueness makes them equal everywhere. Hence $z_0\in\Delta(X)$. The diagonal image is stable under specialization and is closed by Lemma 4.5; Proposition 4.2 gives separatedness.

For properness, the existence-and-uniqueness hypothesis first gives separatedness. After a noetherian base change $Y'\to Y$, take a closed subscheme $Z\subseteq X\times_YY'$, a point $z_1\in Z$, and a proper specialization $f'(z_1)=y_1\leadsto y_0$. The extension $k(z_1)/k(y_1)$ is finitely generated. Apply (a) to the local ring of $y_0$ in the reduced closure of $y_1$ to obtain a DVR of $k(z_1)$ dominating it. The existence hypothesis gives a lift to $X$, hence to $X\times_YY'$. Its generic point lies in $Z$, so the lift factors through $Z$, and its closed point lies over $y_0$. Thus $f'(Z)$ is stable under specialization and closed by Lemma 4.5.

To pass to an arbitrary affine base change $\operatorname{Spec}B\to\operatorname{Spec}A\subseteq Y$, let $Z\subseteq X_B$ be closed and let $b\notin f_B(Z)$. Choose a quasi-compact open $V\subseteq X_B\setminus Z$ containing the fiber over $b$. Since $X$ is of finite presentation over $A$, the finitely many equations describing $V$ descend to a finite-type $A$-subalgebra $B_i\subseteq B$. Enlarging $B_i$ makes the fiber over the image $b_i$ lie in the descended open $V_i$: emptiness of the complementary fiber is witnessed on finitely many affine charts by finitely many equations and denominators. The image of $X_{B_i}\setminus V_i$ is closed by the noetherian case and avoids $b_i$. Its complement pulls back to a neighborhood of $b$ disjoint from $f_B(Z)$. Thus every base change is closed, and $f$ is proper.
